\documentclass[10pt,a4paper]{article}
\usepackage[utf8]{inputenc}
\usepackage[T1]{fontenc}
\usepackage{lmodern}
\usepackage[margin=1in]{geometry}
\usepackage{xcolor}
\usepackage{amsmath,amssymb,amsfonts}
\usepackage{graphicx}
\usepackage{booktabs}
\usepackage{adjustbox}
\usepackage{tcolorbox}
\usepackage{enumitem}
\usepackage{microtype}
\usepackage{array}
\usepackage{fancyhdr}

% Custom Colors matching house style
\definecolor{bitsnavy}{RGB}{33,53,94}
\definecolor{bitsblue}{RGB}{44,90,160}
\definecolor{bitsgreen}{RGB}{46,125,82}
\definecolor{bitsamber}{RGB}{200,136,30}
\definecolor{bitsred}{RGB}{178,58,72}

% Callout Boxes
\tcolorboxenvironment{intuition}{colback=bitsblue!8,colframe=bitsblue,title=\textbf{Everyday Picture \& Intuition},fonttitle=\bfseries,boxrule=0.8pt,arc=3pt}
\tcolorboxenvironment{worked}{colback=bitsgreen!8,colframe=bitsgreen,title=\textbf{Fully Worked Example},fonttitle=\bfseries,boxrule=0.8pt,arc=3pt}
\tcolorboxenvironment{everyday}{colback=bitsnavy!5,colframe=bitsnavy,title=\textbf{Real-World Picture},fonttitle=\bfseries,boxrule=0.8pt,arc=3pt}
\tcolorboxenvironment{watchout}{colback=bitsred!8,colframe=bitsred,title=\textbf{Watch Out! Pitfalls},fonttitle=\bfseries,boxrule=0.8pt,arc=3pt}
\tcolorboxenvironment{keytake}{colback=bitsamber!10,colframe=bitsamber,title=\textbf{Key Takeaways},fonttitle=\bfseries,boxrule=0.8pt,arc=3pt}

\newenvironment{intuition}{\begin{tcolorbox}}{\end{tcolorbox}}
\newenvironment{worked}{\begin{tcolorbox}}{\end{tcolorbox}}
\newenvironment{everyday}{\begin{tcolorbox}}{\end{tcolorbox}}
\newenvironment{watchout}{\begin{tcolorbox}}{\end{tcolorbox}}
\newenvironment{keytake}{\begin{tcolorbox}}{\end{tcolorbox}}

\newenvironment{steps}{\begin{enumerate}[label=\textbf{Step \arabic*:},leftmargin=*]}{\end{enumerate}}

\newcommand{\secsub}[1]{\noindent\textit{\large #1}\vspace{0.8em}}
\newcommand{\housefig}[2]{
\begin{figure}[htbp]
\centering
\includegraphics[width=0.85\linewidth]{#1}
\caption{#2}
\end{figure}
}

\pagestyle{fancy}
\fancyhf{}
\rhead{\color{bitsnavy}\bfseries WILP, BITS Pilani}
\lhead{\color{bitsnavy}\bfseries ISM Session 12 Study Companion}
\rfoot{Page \thepage}

\begin{document}

% Banner
\begin{tcolorbox}[colback=bitsnavy,colframe=bitsnavy,arc=0pt,outer arc=0pt,top=12pt,bottom=12pt]
\color{white}
\centering
{\large\bfseries WILP, BITS Pilani}\\[4pt]
{\LARGE\bfseries Session 12: Multiple Regression \& Intro to Time Series}\\[4pt]
{\small Course: Introduction to Statistical Methods (ISM) \quad|\quad Instructor: Dr. YVK Ravi Kumar}
\end{tcolorbox}

\vspace{1em}

\section{Multiple \& Non-Linear Regression}
\secsub{Extend simple lines to multiple predictors and curved non-linear patterns.}

Multiple linear regression models a continuous target variable $Y$ using multiple explanatory variables: $Y = \beta_0 + \beta_1 X_1 + \beta_2 X_2 + \dots + \beta_k X_k + \epsilon$. Non-linear models apply transformations like logarithms or quadratic terms when data follows curves.

\begin{intuition}
Predicting house price is not just about square footage. You also need bedrooms, neighborhood age, and distance to public transport. Multiple regression assigns weights to all predictors simultaneously.
\end{intuition}

\subsection{Non-Linear Transformations}
For diminishing returns, logarithmic transformations $Y = \beta_0 + \beta_1 \ln(X)$ or quadratic forms $Y = \beta_0 + \beta_1 X + \beta_2 X^2$ fit non-linear curved trends while remaining linear in parameters.

\housefig{figures/nonlinear_regression.png}{Polynomial and logarithmic regression curves fitting non-linear relationships.}

\begin{worked}
Consider fitting a quadratic regression $Y = 5.0 + 2.0X - 0.1X^2$ to estimate yield ($Y$) based on fertilizer dosage ($X$).
\begin{steps}
\item Evaluate predicted yield at $X = 5$:
\begin{equation}
\hat{Y} = 5.0 + 2.0(5) - 0.1(5^2) = 5.0 + 10.0 - 2.5 = 12.5
\end{equation}
\item Evaluate predicted yield at $X = 10$:
\begin{equation}
\hat{Y} = 5.0 + 2.0(10) - 0.1(10^2) = 5.0 + 20.0 - 10.0 = 15.0
\end{equation}
Notice how doubling $X$ from 5 to 10 does not double the yield due to diminishing returns.
\end{steps}
\end{worked}

\section{Time Series Components}
\secsub{Decompose historical sequence observations into structural components.}

A time series is a set of quantitative observations ordered over successive periods. It breaks down into four primary components: Secular Trend ($T_t$), Seasonal variation ($S_t$), Cyclical variation ($C_t$), and Irregular noise ($I_t$).

\begin{intuition}
Think of sales data like ocean waves. The tide rising over years is the trend. Summer beach surges are seasonality. Multi-year economic booms are cycles. Random unexpected storm days are irregular noise.
\end{intuition}

\subsection{Additive vs. Multiplicative Models}
\begin{itemize}
\item \textbf{Additive Model:} $y_t = T_t + S_t + C_t + I_t$ (used when seasonal swings remain constant in dollar magnitude).
\item \textbf{Multiplicative Model:} $y_t = T_t \times S_t \times C_t \times I_t$ (used when seasonal swings grow proportionally with overall trend level).
\end{itemize}

\housefig{figures/time_series_components.png}{Decomposition of time series into trend, seasonal, cyclical, and irregular noise components.}

\section{Stationarity and the Split-Half Test}
\secsub{Check whether time series statistical properties remain constant over time.}

A time series is stationary if its mean, variance, and autocorrelation structure do not depend on time. Non-stationary series must be transformed (e.g., via first differencing $w_t = y_t - y_{t-1}$) prior to modeling.

\begin{intuition}
A stationary process is like a steady thermostat reading around 72 degrees. A non-stationary process is like a room heating up over time where the average temperature keeps rising.
\end{intuition}

\housefig{figures/stationarity_split.png}{Split-half test comparing sub-sample means to diagnose non-stationarity.}

\begin{worked}
Given time series data ($n = 10$): $y_t = [4, 6, 5, 7, 9, 8, 10, 12, 11, 13]$.
\begin{steps}
\item Split data into two equal halves:
First half ($t=1..5$): $4, 6, 5, 7, 9 \implies \bar{y}_1 = \frac{31}{5} = 6.20$.
Second half ($t=6..10$): $8, 10, 12, 11, 13 \implies \bar{y}_2 = \frac{54}{5} = 10.80$.
\item Compare means: Means differ significantly ($10.80 - 6.20 = 4.60$), indicating non-stationarity due to upward trend.
\item Apply first differencing $w_t = y_t - y_{t-1}$:
$w_t = [6-4, 5-6, 7-5, 9-7, 8-9, 10-8, 12-10, 11-12, 13-11] = [2, -1, 2, 2, -1, 2, 2, -1, 2]$.
New differenced mean $\bar{w} \approx 1.0$, which remains stable over time. Differencing once ($d=1$) restores stationarity!
\end{steps}
\end{worked}

\begin{watchout}
Fitting standard linear regression directly on non-stationary time series leads to spurious regression with inflated $R^2$ values.
\end{watchout}

\begin{keytake}
\begin{itemize}
\item Multiple regression isolates simultaneous impacts of several predictor variables.
\item Time series data decomposes into Trend, Seasonal, Cyclical, and Irregular components.
\item Non-stationary time series have changing means or variances over time and require differencing.
\end{itemize}
\end{keytake}

\end{document}
